\section*{Introduction}
\phantomsection
\addcontentsline{toc}{section}{Introduction}
\label{v:presentacion}
This work derives the Bose--Einstein and Fermi--Dirac distributions
from quantum completions of states carrying trajectory and memory,
and proves their convergence to the dilute Maxwell--Boltzmann regime.
The same bosonic statistics, realized on the transverse modes of the
electromagnetic field, determines the Planck spectral law, the
Stefan--Boltzmann flux, and the Wien displacement laws. The central
problem is to relate state structure, admissible occupations, and thermal
observables through explicit operators and measures. Occupation and
radiation laws are thus obtained at connected stages of a construction,
with the equilibrium and representation conditions appropriate to each.

The foundation is Triadic Modular Holography (HMT), together with its
dimensional realization in Dimensional Mechanics. The \emph{All Possible
Paths} (APP) construction is organized on the support
\((\mathbb Z/9\mathbb Z)^2\). Translations by
\(\pm(1,0)\) and \(\pm(0,1)\) define its additive Cayley graph;
the square cells provide a two-dimensional toroidal structure.
Addition and multiplication are evaluated on this same support while
retaining remainders and quotients. The oriented ternary signature,
TRIT, takes values in \(\{-1,0,+1\}\) and distinguishes orientation and
regime. Its integer lift retains both the residue and the carry.
The \emph{Topological Phase Kernel} (TPK)
composes phase selection, transport, update, and state prolongation.
Its record retains trajectory and memory: phase return does not reset
the state.

Initial conditions and compatible refinement produce the genealogical
sections used in the linear realization. Regional generation, the
arithmetic of histories, the twelve-phase record, the determination of
alpha, and exceptional incidence establish the antecedents of the
coefficients and their representations. The state space, unitary
evolution, and multiparticle occupation are constructed on this domain.
Compatibility with depth restrictions precedes transport between
realizations; in the finite instance, the calendar provides an explicit
Hamiltonian, propagator, and action functional.

The decisive step toward statistics is the compositional parity
character. The even sector is completed through symmetric powers, with
nonnegative integer occupations; the odd sector is completed through
exterior powers, whose creation operators are nilpotent and whose
occupation operators are projections. This yields algebraic Pauli
exclusion. The relativistic identification of this parity with spin
additionally retains rotation-representation compatibility and locality
conditions. The multiplicities determine the microcanonical counts,
and factorization of the Gibbs traces yields the Bose--Einstein and
Fermi--Dirac distributions. In the thermodynamic regime, variation of
the combinatorial entropy under energy and population constraints provides a second
characterization of equilibrium. The bosonic sector retains the
chemical-potential condition ensuring convergence. Dilute-regime bounds
quantify the common Maxwell--Boltzmann limit.

The radiative realization applies bosonic occupation to a
three-dimensional periodic cavity with two transverse polarizations,
linear dispersion, and zero chemical potential. The constant mode is
fixed as background data before the excitation state is formed.
Mode density and energy per excitation produce the Planck spectrum.
Convergence from the discrete sums is proved by uniformly controlling
the neighborhood of zero frequency and the high-frequency tail.
Spectral moments give number, energy, and entropy densities; angular
integration gives the Stefan--Boltzmann flux. Conversion between
frequency and wavelength also transforms the measure through its
Jacobian, so the two Wien maxima satisfy different equations.
Each displacement constant is characterized by the maximum of the
corresponding density.

The relationship between action, calendar, and temperature determines
the dimensional normalizations of this realization. Thermal transduction
distinguishes decision energy, dimensionless information, and the
Boltzmann--temperature pair. Its composition with constitutive speed
connects the temporal, energetic, and spatial scales of radiation.
Continuity, generator-domain, contractivity, and convergence assumptions
remain attached to the operations that require them.

Between transduction and the completions, the carrier-fixing chain consists of
tritic renewal, double grading, positive energy memory, Williamson spectrum,
threefold refinement, marked section, boundary currents, and centred thermal
equivalence. This chain retains the origin that disappears from a
normalised density and transports it to area, free energy, gravitation, and
radiation. The complete evaluation is therefore presented not as a decimal
coincidence, but as the output of typed operators, fibres, and maps.

\Needspace{13\baselineskip}
The entropic interpretation distinguishes occupation configurations,
the measure on histories, and the spectrum of a density operator.
A calendar's chronological frequency is not the stationary measure
on its envelope of histories. The latter is characterized variationally
through a relative-divergence identity. In the number-conserving,
gauge-invariant fermionic quasifree sector, the one-body correlator
\(0<C<I\) and its modular generator \(K_1\) determine each other in
finite dimension. Transport, reduction, and conditioning specify which
information remains observable. Record recovery and orbital
identification of operators are linked through a bilateral balance
that preserves state and memory together
(Section~\ref{sec:acoplamiento-recuperacion}).

The areal unit and entanglement complete this relationship between
structure and information. The finite modular identity expresses
entropy variations through sectoral areas minus a nonnegative
relative-entropy deficit; for \(\Delta_{\rm mod}\ne0\) and known
normalization, area and the modular Hamiltonian recover the occupations
of both sectors. The connection between memory, occupation, and radiation thus
retains the maps connecting successive objects as well as their scalar
values. Appendix~\ref{v:integracion} specifies the realization domains
and the conditions of these compositions.

The finite realization on \(\Lambda=(\mathbb Z/27\mathbb Z)^2\)
establishes exact support and entropy bounds, with saturations by
subgroups and annihilators. On the dodecaphasic domain, \(\nu_{120}\) and
\(\nu_{270}\) detect ordered correlations erased by the tritic histogram.
On the prime domain, \(R_{\rm vac}\) composes the generated Sturmian word,
the internal realization of irreducibles, and the \(\log p\) clock. Its
modal identity is understood through symmetric sums or Fejér means, and
its tables constitute finite numerical witnesses.

\subsection*{Relations between the constructions and their results}
The following table identifies the contribution of each block to the
argument. The exposition preserves the common antecedents and then
distinguishes the specific operations: representation, occupation,
conditioning, and radiative reading. The graded representation determines occupation; the spatial realization
and dispersion determine its radiative application.

\begin{table}[H]
\centering
\small
\renewcommand{\arraystretch}{1.2}
\begin{tabularx}{\textwidth}{@{}>{\raggedright\arraybackslash}p{.29\textwidth}>{\raggedright\arraybackslash}X@{}}
\toprule
Construction and location & Construction and role in the argument \\
\midrule
Core and extensions, §§\ref{sec:nucleo}--\ref{exc:seccion}
& Admissible states, oriented routes, memory, and refinement; regional
outputs, dodecaphase record, and incidence. These antecedents preserve
the origin of the coefficients and subsequent representations. \\
Quantum realization and action,
§§\ref{v:rec:sec:postulados-cuanticos-genealogicos-rev2}--\ref{sec:hmtf2-cristal-temporal}
& State space, unitary evolution, and calendar action. The realization
fixes the operators that subsequently enter the energies and
equilibrium states. \\
Memory and temperature,
§\ref{v:ant:sec:ii-kb-aut-desarrollo}
& Readout fibers, reversibility, and the Boltzmann--temperature pair.
Dimensionless information, energy scale, and thermal chart are distinguished. \\
Spectral capacity and marked bridge,
§\ref{v:sec:capacidad-espectral-termica}
& Weighted microscopic tree, graded carrier, boundary currents, capacity and
Williamson scales, and the variational thermometric transducer. \\
Parity, Fock completions, and equilibrium distributions,
§\ref{v:sec:completaciones-algebraicas} and
§§\ref{v:rec:fock_gibbs:section:01}--\ref{v:rec:ocupacion:section:03};
milestones \ref{v:rec:completaciones:statement:02},
\ref{v:rec:fock_gibbs:statement:01},
\ref{v:rec:thm:factorizacion-gran-canonica-rev3}, and
\ref{v:rec:thm:limite-clasico-fd-be-rev2}
& Parity, Fock space, refinement transport, and the Gibbs state.
Multiplicities give the counts; the traces yield the two occupation
distributions, whose dilute limit is quantified, §§\ref{v:rec:ocupacion:section:02}--\ref{v:rec:ocupacion:section:03}. \\
Histories, area, and entanglement
& Areal--volumetric incidence determines the history measure;
state reduction and the area operator support the modular reading
of entropy. See Theorem~\ref{thm:incidencia-hojas-anexo}
and Section~\ref{v:ant:sec:area}. \\
Fourier uncertainty, ordered memory, and prime vacancies,
§\ref{v:n19:sec:incertidumbre-fourier} and
§\ref{sec:f051-prime-vacancy-observer}
& The finite realization establishes the support and entropy bounds;
dodecaphase readers distinguish order from marginals; the prime-vacancy
observer resolves the correlation through the symmetric/Fejér modal
identity, with a historical witness through \(10^8\), independent
reproduction through \(10^6\), and a critical bound conditional on modal
uniformity. \\
Bosonic occupation and thermal radiation laws,
§\ref{v:sec:radiacion-termica}; Propositions
\ref{v:prop:ley-espectral}, \ref{v:prop:stefan}, and \ref{v:prop:wien};
limit in Lemma~\ref{v:lem:limite-termodinamico}
& Transverse modes, constitutive speed, and bosonic statistics determine
the thermal sums. Estimates near the origin and in the tail justify
the limit; its moments yield densities, flux, and spectral maxima. \\
\bottomrule
\end{tabularx}
\caption{Antecedent constructions and their subsequent use in the
statistical and radiative argument.}
\label{v:tab:mapa-demostrativo}
\end{table}

The presence of shared antecedents reflects this actual dependence.
Each block preserves the information required by the next operation;
a coincident dimension or scalar value is accompanied by the map that
allows its reuse. Statistics thus receives a graded representation and
a Hamiltonian, while radiation additionally receives a spatial
realization and its dispersion relation.
